\section{Conclusion}
\label{sec:conclusion}

A player's first randomly assigned tournament opportunity produces larger career gains than subsequent ones, and these gains operate through access rather than ability and dissipate within a year. In the first-time ATP Grand Slam sample ($N = 120$), lottery selection raises ranking points by 72.5 at 4 weeks ($p = 0.008$), 82.8 at 12 weeks ($p = 0.002$), and 72.6 at 26 weeks (clustered $p = 0.047$; Fisher $p = 0.164$), fading by 52 weeks as the rolling ranking window absorbs the initial gain. These point estimates are roughly twice the effects estimated on the full sample that pools first-time and repeat LL recipients, and a mechanical decomposition into stratum-weighted components shows that $\hat\beta^{(n=0)} / \hat\beta^{\text{pool}} \approx 2$ across horizons on the ATP side. On the WTA tour ($N = 82$), effects concentrate at 8 weeks ($+55.0$, $p = 0.017$), and Elo gains at medium horizons align with a genuine ability channel on the women's tour. Non-Grand-Slam control-function estimates on larger samples ($N = 2{,}054$ ATP; $N = 1{,}423$ WTA) reproduce the ATP access pattern: ranking-point gains of 20.6 at 4 weeks ($p = 0.041$) and 44.6 at 26 weeks ($p = 0.002$). Wald tests reject equality of first-LL and second-LL effects for tournament access ($p = 0.002$ ATP; $p = 0.009$ WTA in the non-GS sample), and third-plus LL recipients show weak or negative effects.

The dominant mechanism is access accumulation. The ranking system does not adjust for opponent quality: a first-round win at a Grand Slam earns the same points regardless of the opponent's strength. When a first-time lucky loser earns points at the focal event, the improved ranking grants entry to subsequent tournaments and further chances to accumulate points. The match-level tournament performance model reports $\hat{\delta} = -0.141$ (player-clustered $p = 0.17$) for first-time ATP GS recipients: directionally consistent with overseeding when novice entrants face opponents above their ability level, though the small first-time GS sample does not identify the effect from zero. Ranking-point gains therefore come from more draws entered and more matches played rather than from a per-match performance improvement. The WTA evidence differs modestly, with Elo gains at medium horizons pointing to a first-main-draw learning channel that is absent from the pooled-LL results.

The ranking feedback loop, in which points determine access and access determines point-earning opportunities, amplifies a one-time shock into a multi-month cascade; the cascade is self-limiting because sustained ability dominates a single favorable shock within a year when performance is continuously measured and positions are regularly updated. The comparison with \citet{Oyer2006_initial_placement} and \citet{Oreopoulos2012_graduating_recession}, where initial conditions cause decade-long career effects in less transparent markets, points to market transparency as the key moderator: the 52-week rolling window provides a natural correction mechanism that other labor markets lack.

Two limitations remain. The Grand Slam lottery sample is small ($N = 120$ ATP, $N = 82$ WTA), limiting power for heterogeneity analysis. Fisher randomization inference, more appropriate than asymptotic inference for this small-sample lottery, supports the main effects at 4--12 weeks but not at 26 weeks, so the middle horizon should be read with appropriate caution. Lottery contamination is an additional caveat: an unknown fraction of observations were assigned by ranking rather than by draw, and ranking-rule assignment selects higher-ranked players. The pre-treatment controls absorb the observable component of this bias, within-event fixed effects restrict comparisons to narrow ranking gaps, and a verified-lottery subsample built from a data-only criterion (rank-among-losers greater than one) recovers a 26-week ATP ranking-point effect of 76.6 ($p = 0.050$) and an 8-week WTA effect of 77.5 ($p = 0.005$) on 35 ATP and 17 WTA verified events, in line with the headline estimates (Section~\ref{sec:robustness_verified}). A full resolution of the remaining 23 ambiguous entries would require withdrawal-timing data not available in the public record. Future work could measure longer-run outcomes and test whether the first-break phenomenon generalizes to other tournament-style markets with threshold-based access.
